\section{单群与交错群的单性}

若群$G$有非平凡真正规子群$N$，就可以同时研究$N$及商群$G/N$，以便逐步了解$G$。不能再这样分解的群称为单群，它们是研究一般群结构的重要组成部分。

\begin{definition}
非平凡群$G$若没有除$\{e\}$与$G$以外的正规子群，则称为\textbf{单群}。
\end{definition}

\begin{proposition}
非平凡交换群为单群，当且仅当它是素数阶循环群。
\end{proposition}
\begin{proof}
交换群的全部子群都正规，故只需判定何时没有非平凡真子群。若$G$交换且单，取$a\ne e$，则$\langle a\rangle$非平凡且正规，只能等于$G$，所以$G$循环。

若$G$无限，则$\langle a^2\rangle$为非平凡真子群；若其阶$m$为合数，则取$1<d<m$且$d\mid m$，由循环群子群分类得到$d$阶非平凡真子群。两种情况都与单性矛盾，故$|G|$只能为素数。
反过来，素数阶群由拉格朗日定理没有非平凡真子群，因此当然是单群。
\end{proof}

非交换单群的情形复杂得多。下面给出一类非交换单群：$n\geq5$时的交错群$A_n$。为此先研究对换、3-轮换以及共轭对轮换分解的影响。

\begin{lemma}\label{p2:three-cycles}
每个置换都可表示为对换的乘积；每个偶置换都可表示为3-轮换的乘积。因此$A_n$由全体3-轮换生成。
\end{lemma}
\begin{proof}
将置换分解为不相交轮换后，对每个轮换使用
\[
(i_1i_2\cdots i_r)=(i_1i_r)\cdots(i_1i_2).
\]
右侧按从右到左作用，把$i_1$送到$i_2$，把$i_k$送到$i_{k+1}$（$2\leq k<r$），把$i_r$送回$i_1$，并固定其余点，所以等式成立。这就得到对换分解。

每个对换为奇置换，故偶置换的任一对换分解都含偶数个对换。将它们按次序两两配对。若一对中的两个对换相同，乘积为$e$；若不同，则要么恰有一个共同点，要么没有共同点。重新命名涉及的点后，两种乘积分别为
\begin{align*}
(ij)(ik)&=(ikj),\\
(ij)(kl)&=(ij)(jk)(jk)(kl)\\
&=(jki)(klj).
\end{align*}
第一式将$i,k,j$依次送到$k,j,i$；第二式在中间插入$(jk)^2=e$，再把左右两对分别化为3-轮换。因此每一对都能写为3-轮换的乘积，任意偶置换也能如此表示。每个3-轮换又是偶置换，故它们生成的正是$A_n$。
\end{proof}

\begin{proposition}
共轭把轮换中的各个符号逐个替换：
\[
\pi(i_1\cdots i_r)\pi^{-1}
=(\pi(i_1)\cdots\pi(i_r)).
\]
若在不相交轮换分解中计入长度为1的轮换，则称各个轮换长度（保留重复次数、不计次序）为置换的\textbf{轮换类型}。两个置换在$S_n$中共轭，当且仅当其轮换类型相同。若$n\geq5$，全部3-轮换在$A_n$中也互相共轭。
\end{proposition}
\begin{proof}
记$\pi(i_k)=j_k$。对$1\leq k<r$，
\[
\pi(i_1\cdots i_r)\pi^{-1}(j_k)
=\pi(i_{k+1})=j_{k+1}.
\]
同理$j_r$被送到$j_1$。若点不在$\{j_1,\ldots,j_r\}$中，它的$\pi^{-1}$像也不在原轮换中，故共轭后仍固定。这就证明公式。

把公式用于不相交轮换分解，便知共轭保持轮换类型。反过来，若两个置换轮换类型相同，则将等长轮换配对，并依循环次序把每个点映到配对轮换中的对应点；不动点之间也任选一一对应。这定义一个$\pi\in S_n$，由公式知它把第一个置换共轭到第二个。

最后，取任意两个3-轮换$(i_1i_2i_3)$与$(j_1j_2j_3)$，先选$\pi\in S_n$满足$\pi(i_k)=j_k$，$k=1,2,3$。若$\pi$为偶置换，已满足要求。若$\pi$为奇置换，由$n\geq5$可在$i_1,i_2,i_3$之外取两个不同的点$l,m$，令$\varphi=\pi(lm)$。此时$\varphi$为偶置换，且仍有$\varphi(i_k)=j_k$，所以
\[
\varphi(i_1i_2i_3)\varphi^{-1}=(j_1j_2j_3).
\]
故全部3-轮换在$A_n$内共轭。
\end{proof}

\begin{example}
若$\sigma=(123)(45)$，$\pi$把$1,2,3,4,5$分别映到$5,4,1,3,2$，则只需按上述公式替换符号，就得到
\[
\pi\sigma\pi^{-1}=(541)(32).
\]
\end{example}

